\section*{الفصل \ref{ch:b1:stereochemistry-in-depth} --- الكيمياء الفراغية بتعمق}

\begin{solution}{exo:b1:stereochemistry-in-depth:1}
\omterm{def:b1:stereochemistry-in-depth:isomers}{متماكبان بنيويان}؛ \omterm{def:b1:stereochemistry-in-depth:enantiomers}{متماكبان مرآويان}؛ \omterm{def:b1:stereochemistry-in-depth:enantiomers}{متماكبان لامرآويان} (\omterm{def:b1:stereochemistry-in-depth:isomers}{متماكبان فراغيان} ليس أحدهما
صورة مرآوية للآخر)؛ تشكّلان لمركب واحد؛ \omterm{def:b1:stereochemistry-in-depth:enantiomers}{متماكبان لامرآويان}
(\cip{R,S} هو شكل الميزو).
\end{solution}

\begin{solution}{exo:b1:stereochemistry-in-depth:2}
\babelsublr{\ce{-OH} $>$ \ce{-NH2} $>$ \ce{-CH3} $>$ \ce{-H} (O, N, C, H)}.
\babelsublr{\ce{-COOH} (O,O,O) $>$ \ce{-CHO} (O,O,H) $>$ \ce{-CH2OH} (O,H,H) $>$
\ce{-CH(CH3)2} (C,C,H)}. \babelsublr{\ce{-Br} $>$ \ce{-Cl} $>$ \ce{-CCl3} (Cl,Cl,Cl)
$>$ \ce{-CH2Cl} (Cl,H,H)}: فالذرة الأولى تحسم قبل جاراتها.
\end{solution}

\begin{solution}{exo:b1:stereochemistry-in-depth:3}
2,3-ثنائي كلوروبوتان: مركزان متكافئان، \cip{R,R} و\cip{S,S}
وشكل الميزو \cip{R,S}: ثلاثة. 2-برومو-3-كلوروبوتان: مركزان مختلفان،
أربعة \omterm{def:b1:stereochemistry-in-depth:isomers}{متماكبات فراغية}، ولا شكل ميزو. البنتان-2,4-ثنائي أول: ثلاثة، و\cip{2R,4S}
هو الميزو.
\end{solution}

\begin{solution}{exo:b1:stereochemistry-in-depth:4}
على C5: \babelsublr{\foreignlanguage{arabic}{الإيزوبروبينيل} $>$ C6 $>$ C4
$>$ H}، فللإيزوبروبينيل (C,C,C) إذ تُعدّ الرابطة الثنائية مرتين. ويحمل C6 وC4 كلاهما (C,H,H)؛ وبعد خطوة أخرى يقود C6 إلى
كربون الكربونيل (O,O,C) ويقود C4 إلى C3 (C,C,H): فيتقدم C6. ومع مجموعة الإيزوبروبينيل
في الأمام وH في الخلف، يدور الترتيب \babelsublr{\foreignlanguage{arabic}{إيزوبروبينيل} $\to$ C6 $\to$ C4}
في اتجاه عقارب الساعة كما هو مرسوم: \cip{R}، كارفون النعناع؛ وذو الإسفين المخطط
\cip{S}.
\end{solution}

\begin{solution}{exo:b1:stereochemistry-in-depth:5}
$\theta = 0$: الميثيلان متطابقان، \qty{16.5}{kJ/mol}؛ $60^\circ$: منحرف، نحو
\qty{2.9}{kJ/mol} (الحد الأدنى 2.8 عند $62^\circ$)؛ $120^\circ$: ميثيل يطابق H،
\qty{15.2}{kJ/mol}؛ $180^\circ$: مضاد، 0. ومع تساوي الطاقات، \omterm{def:b1:stereochemistry-in-depth:conformations}{تشكّلان
منحرفان} مقابل \omterm{def:b1:stereochemistry-in-depth:conformations}{تشكّل مضاد} واحد: فثلثا الجزيئات منحرفة.
\end{solution}

\begin{solution}{exo:b1:stereochemistry-in-depth:6}
السلسلة رأسية و\ce{COOH} في الأعلى و\ce{CH3} في الأسفل؛ والرتب
$\ce{NH2} > \ce{COOH} > \ce{CH3} > \ce{H}$. وبوضع \ce{NH2} على اليسار
وH على اليمين، يدور $\ce{NH2} \to \ce{COOH} \to \ce{CH3}$ في اتجاه عقارب الساعة
كما يُرى؛ وH أفقية (نحو القارئ)، فيُعكس
الواصف: \cip{S}. فمجموعة الأمين على اليسار.
\end{solution}

\begin{solution}{exo:b1:stereochemistry-in-depth:7}
\textit{cis}: على كربونين متجاورين، رابطة إلى الأعلى وأخرى إلى الأسفل من بين المواضع \omterm{def:b1:stereochemistry-in-depth:conformations}{المحورية}
\omterm{def:b1:stereochemistry-in-depth:conformations}{والاستوائية} تعطي محوري--استوائي في الكرسيين كليهما.
\textit{trans}: ثنائي \omterm{def:b1:stereochemistry-in-depth:conformations}{الاستوائية} أو ثنائي \omterm{def:b1:stereochemistry-in-depth:conformations}{المحورية}. وفي \omterm{def:b1:stereochemistry-in-depth:conformations}{الكرسي} الثنائي \omterm{def:b1:stereochemistry-in-depth:conformations}{الاستوائية} يكون
الميثيلان منحرفين أحدهما بالنسبة إلى الآخر (تآثر واحد)؛ وفي المتماكب \textit{cis}
هذا التآثر إضافة إلى ميثيل \omterm{def:b1:stereochemistry-in-depth:conformations}{محوري}، أي نحو \qty{5.7}{kJ/mol}
أكثر. فالمتماكب \textit{trans} هو الأكثر استقرارًا.
\end{solution}

\begin{solution}{exo:b1:stereochemistry-in-depth:8}
$[\alpha] = 13.2/(2.00 \times 0.100) = +66.0^\circ$. وأنبوب طوله \qty{1.00}{dm}:
$+6.60^\circ$. \omterm{def:b1:stereochemistry-in-depth:enantiomers}{والمتماكب المرآوي}: $-13.2^\circ$ في الأنبوب ذي الطول \qty{2.00}{dm}.
\end{solution}

\begin{solution}{exo:b1:stereochemistry-in-depth:9}
\qty{25}{\celsius}: $\exp(5700/(8.314 \times 298.15)) = 10.0$، والكسر
$10.0/11.0 = \qty{91}{\percent}$. \qty{-50}{\celsius}: $\exp(5700/(8.314
\times 223.15)) = 21.6$، أي \qty{96}{\percent}.
\end{solution}

\begin{solution}{exo:b1:stereochemistry-in-depth:10}
في الإيثان ذرات الهيدروجين الست كلها متماثلة: فكل \omterm{def:b1:stereochemistry-in-depth:conformations}{تشكّل متخالف}
هو \omterm{def:b1:stereochemistry-in-depth:isomers}{التشكّل} نفسه. أما في البوتان فيمكن أن يقف الميثيل الخلفي مقابل الأمامي
(مضاد) أو بجانبه ($\pm 60^\circ$، منحرف)، حيث يزاحم الميثيلان
أحدهما الآخر. الأعداد النسبية: مضاد 1، ومنحرف $2\exp(-2.8/2.48) = 0.65$: أي نحو
\qty{61}{\percent} مضاد و\qty{39}{\percent} منحرف.
\end{solution}

\begin{solution}{exo:b1:stereochemistry-in-depth:11}
C2 وC4 مركزان فراغيان. ويحمل C3 كلًّا من H وOH وفرعين
متماثلين في البنية؛ ولا يختلفان إلا حين يحمل الكربونان C2 وC4 واصفين
متعاكسان، وعندئذ فقط يكون C3 مركزًا فراغيًا (شبه لامتناظر). المتماكبات:
\cip{2R,4R} و\cip{2S,4S}، زوج من \omterm{def:b1:stereochemistry-in-depth:enantiomers}{المتماكبات المرآوية} (C3 ليس مركزًا فراغيًا)؛
و\cip{2R,4S} مع الترتيبين عند C3، وهما شكلا ميزو. أي أربعة إجمالًا.
\end{solution}

\begin{solution}{exo:b1:stereochemistry-in-depth:12}
$[\alpha]_{\text{العيّنة}} = 2.31/(1.00 \times 0.050) = +46.2^\circ$،
و$2x - 1 = 46.2/66.0 = 0.700$: $x = 0.850$، أي \qty{85}{\percent} من
\omterm{def:b1:stereochemistry-in-depth:enantiomers}{المتماكب المرآوي} $(+)$ و\qty{15}{\percent} من $(-)$.
\end{solution}

\begin{solution}{pb:b1:stereochemistry-in-depth:1}
\textbf{1.} حلقة الهكسان الحلقي: يحمل C1 المجموعة OH، وC2 مجموعة الإيزوبروبيل، وC5
الميثيل؛ وC1 وC2 وC5 مراكز فراغية.
\textbf{2.} $2^3 = 8$، أي أربعة أزواج من \omterm{def:b1:stereochemistry-in-depth:enantiomers}{المتماكبات المرآوية}.
\textbf{3.} لا ترتيب فيه مستوي مرآة: فالمتبادلات الثلاثة
مختلفة كلها.
\textbf{4.} \babelsublr{\ce{OH} $>$ C2 (C,C,H) $>$ C6 (C,H,H) $>$ H}.
\textbf{5.} \babelsublr{C1 (O,C,H) $>$ \foreignlanguage{arabic}{الإيزوبروبيل} (C,C,H) $>$ C3 (C,H,H) $>$ H}.
\textbf{6.} يحمل C6 وC4 كلاهما (C,H,H)؛ وبعد ذلك يقود C6 إلى C1 (O,C,H)، وC4 إلى C3
(C,H,H): \babelsublr{C6 $>$ C4 $>$ \ce{CH3} $>$ H}.
\textbf{7.} \cip{1S,2R,5S}.
\textbf{8.} في \omterm{def:b1:stereochemistry-in-depth:conformations}{الكرسي}، تتجه رابطتان استوائيتان على كربونين متجاورين (C1، C2)
إحداهما إلى الأعلى والأخرى إلى الأسفل: \textit{trans}؛ وعلى الكربونين 1 و3 من الحلقة
(C1 وC5، عبر C6) تتجه الرابطتان الاستوائيتان في الجهة نفسها:
\textit{cis}. وهذا هو ترتيب المنثول.
\textbf{9.} تصير الثلاثة كلها \omterm{def:b1:stereochemistry-in-depth:conformations}{محورية}.
\textbf{10.} في \omterm{def:b1:stereochemistry-in-depth:conformations}{الكرسي} المقلوب يكلّف الميثيل وحده نحو
\qty{5.7}{kJ/mol} ومجموعة الإيزوبروبيل الأكبر أكثر؛ وتضيف OH كلفة
أصغر خاصة بها، وكان OH والميثيل المحوريان، على الكربونين 1 و3 وعلى
الوجه نفسه، سيزاحم أحدهما الآخر أيضًا. وهذا أكثر بكثير من \qty{12}{kJ/mol}،
أي نسبة تفوق $\exp(12000/2479) \approx 130$ إلى واحد: فالمنثول كله تقريبًا
في \omterm{def:b1:stereochemistry-in-depth:conformations}{الكرسي} \omterm{def:b1:stereochemistry-in-depth:conformations}{الاستوائي} كليًا.
\textbf{11.} \omterm{def:b1:stereochemistry-in-depth:enantiomers}{متماكب لامرآوي}: لا يختلف إلا مركز واحد من ثلاثة.
\textbf{12.} مع الإيزوبروبيل والميثيل استوائيين، تكون OH \omterm{def:b1:stereochemistry-in-depth:conformations}{محورية}.
\textbf{13.} \omterm{def:b1:stereochemistry-in-depth:enantiomers}{المتماكبات اللامرآوية} تختلف في كل خواصها (هنا
\omterm{def:b1:intermolecular-forces-solvents:hbond}{الرابطة الهيدروجينية} لمجموعة OH \omterm{def:b1:stereochemistry-in-depth:conformations}{محورية} أو \omterm{def:b1:stereochemistry-in-depth:conformations}{استوائية})؛ أما \omterm{def:b1:stereochemistry-in-depth:enantiomers}{المتماكبات المرآوية} فلا تختلف إلا
تجاه الشركاء الكيراليين.
\textbf{14.} $[\alpha] = -2.46/(1.00 \times 0.0500) = -49.2^\circ$، داخل
مجال المرجع.
% ledger: alpha:l-menthol-range
\textbf{15.} $+49.2^\circ$؛ $0$.
\textbf{16.} $-1.97/0.0500 = -39.4^\circ$.
\textbf{17.} $\alpha = \ell c[\alpha]_{(-)}\bigl(x - (1 - x)\bigr) =
\ell c[\alpha]_{(-)}(2x - 1)$.
\textbf{18.} $x = \frac12(1 + \alpha/\alpha_{\text{مرجع}})$ عند $\ell$ و$c$ متساويين.
\textbf{19.} نعم: فكل مركب آخر نشط ضوئيًا يضيف حده الخاص
(\omterm{prop:b1:stereochemistry-in-depth:biot}{قانون بيو}). ويلزم التحقق من النقاوة الكيميائية بطريقة أخرى
(الكروماتوغرافيا) قبل قراءة الدوران على أنه تركيب.
\textbf{20.} $x = \frac12(1 + 1.97/2.46) = 0.900$.
\textbf{21.} $x = \frac12(1 + 2.40/2.46) = 0.988$.
\textbf{22.} الثانية (\qty{98.8}{\percent}).
\textbf{23.} $-3.94^\circ$؛ فيصير خطأ القراءة نفسه كسرًا أصغر
من الزاوية.
\textbf{24.} \textbf{\qty{90.0}{\percent} من المنثول-$(-)$} (و
\qty{10.0}{\percent} من متماكبه المرآوي).
\end{solution}
